\documentclass[../../../include/open-logic-section]{subfiles}

\begin{document}

\olfileid[pt]{his}{set}{mythology}

\olsection{Mais Mito do que História?}

Com mais de um século de perspectiva, ao olhar para estes acontecimentos,
devemos ter cuidado para não tomar estes veredictos por adquiridos. Os resultados eram
certamente novos, entusiasmantes e surpreendentes. Mas quão verdadeiramente chocantes eram?
E demonstraram realmente que não devemos confiar na
intuição geométrica?

Quanto ao choque, \citet{Gouvea2011} observa que a célebre nota de Cantor a Dedekind,
``\emph{je le vois, mais je ne le crois
pas}'' é tirada um tanto fora de contexto. Eis mais desse contexto
(citado de \citeauthor{Gouvea2011}):
\begin{quote}
Peço desculpa pelo meu zelo pelo assunto se lhe faço tantas exigências à
sua bondade e paciência; as comunicações que ultimamente lhe enviei
são mesmo para mim tão inesperadas, tão novas, que não posso ter paz de espírito
até obter de si, amigo honrado, uma decisão sobre a sua
correcção. Enquanto não concordar comigo, só posso dizer:
\emph{je le vois, mais je ne le crois pas.} 
\end{quote}
Cantor sabia que o seu resultado era ``tão inesperado, tão novo''. Mas é
duvidoso que alguma vez achasse o seu resultado \emph{inacreditável}. Como
\citeauthor{Gouvea2011} observa, estava simplesmente a pedir a Dedekind que
verificasse a prova que oferecera. 

\enlargethispage{2\baselineskip}%
Quanto à intuição geométrica: Peano publicou a sua
curva preenchedora de espaço sem incluir quaisquer diagramas. Mas quando Hilbert
publicou a sua curva, explicou o seu propósito: daria aos
leitores um modo claro de compreender o resultado de Peano, se ``se
servirem da seguinte intuição geométrica''; após o que
incluiu uma série de \emph{diagramas} tal como os fornecidos em
\olref[his][set][pathology]{sec}. 

De modo mais geral: embora os diagramas tenham caído um tanto em desuso nas
provas publicadas, não há como escapar ao fato de que os
matemáticos \emph{com frequência} usam diagramas ao provar resultados.
(Grosso modo: bons matemáticos sabem quando podem confiar na
intuição geométrica.)

Em suma: não acredite na \emph{exaltação}; ou, pelo menos, não a aceite por adquirido. Para mais sobre isto, pode ler \citet{Giaquinto2007}.

\end{document}
